\subsection{Constant ring fibers over an entire idempotent semiring}
\label{sec:entire-support}

We now evaluate the amplitude condition when all nonzero coefficient fibers
are one ring and their transition maps are identities. The support semiring
may be infinite and need not be commutative. Its own projective modules
remain part of the classification data.

We prove Theorem~\ref{thm:entire-support-projectives}.



\begin{proof}[Proof of Theorem~\ref{thm:entire-support-projectives}]
First suppose that $M$ is the image of an idempotent endomorphism of
$T^n$. For left modules we use row coordinates: if the image of the
$i$th basis vector has $j$th coefficient $e_{ij}$, the endomorphism is
$x\mapsto xE$, where $(xE)_j=\sum_i x_i e_{ij}$. This convention
keeps the coefficient order correct when the semirings are
noncommutative.

Project $E$ to matrices $E_Q$ over $Q$ and $E_0$ over $R_0$, with
common Boolean support matrix $B$. Let $A$ be the ordinary $R$-matrix
obtained from $E_0$ by replacing absent entries by ring zero. Each
matrix is idempotent, and $A_{ij}=0$ wherever $B_{ij}=0$. The converse
need not hold, since an occupied coefficient can have amplitude zero.
The support of $M$ is
\[
 L=\{l\in Q^n:lE_Q=l\}.
\]
It is a finite free retract. Applying $\B\otimes_Q-$ to the support
retraction identifies $D$ with the image of $B$ in $\B^n$, hence with a
finite projective Boolean semimodule. The finite Boolean classification
therefore makes $D$ a distributive lattice.

For $q\in Q^n$ let $I(q)$ be its set of nonzero coordinates.
Entirety and idempotent addition give
\[
 I(qE_Q)=I(q)B.
\]
Indeed, the $j$th coordinate is nonzero precisely when some
$q_i(E_Q)_{ij}$ is nonzero, equivalently when both of its factors are
nonzero. Thus, for $l\in L$, the map $b$ is $l\mapsto I(l)$ under
the preceding identification. Every fixed subset $I$ of $B$ occurs:
apply $E_Q$ to the vector whose coordinates are $1$ on $I$ and $0$
elsewhere. Its occupied set is $IB=I$.

The free group above $q$ is $R^{I(q)}$, with coordinate inclusions as
transitions. Over a fixed support $l\in L$, the group of $M$ is
\[
 \operatorname{im}\bigl(R^{I(l)}\xrightarrow{\,x\mapsto xA\,}
                              R^{I(l)}\bigr).
\]
This expression depends only on $I(l)$. On the other hand,
$N=\operatorname{im}(E_0)\subseteq R_0^n$ is a finite projective
$R_0$-semimodule with support $D$, and its group at $I\in D$ is
exactly the same image of $A$. Consequently the amplitude diagram of
$M$ is the pullback along $b$ of that of $N$. Nonzero $T$-scalars act
through their ring amplitudes. The classification of projectives over
$R_0$ identifies $N$ with the labelled-poset diagram for finite
projective $R$-modules $V_p$, so it gives the asserted description of
$M$.

For the converse, start with the stated data. The displayed groups and
actions form a $T$-semimodule: transitions are inclusions of summands,
$b(ql)=b(l)$ for $q\ne0$, and the remaining identities are the
ordinary left $R$-module laws. We construct a finite free retraction
explicitly.

Choose $Q$-linear maps
\[
 \sigma:L\longrightarrow Q^n,\qquad
 \rho:Q^n\longrightarrow L,\qquad \rho\sigma=1_L.
\]
Put $l_i=\rho(e_i)$ and $\phi_i=\sigma_i$. Then
$\phi_i(l)l_i\le l$ for every $i,l$. Booleanization gives a split
pair $\sigma_B:D\to\B^n$, $\rho_B:\B^n\to D$, with
\[
 \rho_B(e_i)=b(l_i),\qquad
 (\sigma_Bb(l))_i=\chi(\phi_i(l)).
\]
For every join-irreducible $p\in D$, the equality
\[
 p=\bigvee_{i\in\sigma_B(p)}b(l_i)
\]
supplies a coordinate $i(p)$ with $b(l_{i(p)})=p$ and
$i(p)\in\sigma_B(p)$. For this coordinate and every $d\in D$,
\[
 i(p)\in\sigma_B(d)\quad\Longleftrightarrow\quad p\le d.
\]
The forward implication follows by applying $\rho_B$; the reverse
follows from monotonicity. Hence
\begin{equation}\label{eq:entire-occupancy}
 \phi_{i(p)}(l)\ne0\quad\Longleftrightarrow\quad p\le b(l).
\end{equation}
The coordinate is chosen to be active at $p$, not merely to have
support image $p$.

Choose split $R$-linear maps
\[
 V_p\xrightarrow{u_p}R^{m_p}\xrightarrow{v_p}V_p,\qquad
 v_pu_p=1_{V_p},
\]
with $m_p\ge1$, also for a zero label. Take a free $T$-module with
one block of $m_p$ coordinates for each $p$, followed by the original
$n$ support-frame coordinates. Map the basis elements of the
$p$-block to the values of $v_p$ in the $p$-summand of
$M_{l_{i(p)}}$, and map the final basis elements to the group zeros
at $l_i$. These choices define a $T$-linear map $f:T^N\to M$.
Define a support map $\sigma':L\to Q^N$ by repeating
$\phi_{i(p)}$ in its $p$-block and using $\sigma$ in the final
coordinates. If $\rho'$ is the support map of $f$, then
\[
 \rho'\sigma'(l)=
 \bigvee_p\phi_{i(p)}(l)l_{i(p)}\vee\rho\sigma(l)=l.
\]

On the group at $l$, define a section into the free group over
$\sigma'(l)$ by applying $u_p$ to every summand with $p\le b(l)$
and putting zero amplitudes in the support-frame coordinates.
Equation~\eqref{eq:entire-occupancy} says that precisely the required
$p$-blocks have coefficient group $R$; the other blocks have the
zero group. These maps commute with transitions because both sides
use the same summand inclusions. They commute with a nonzero scalar
$(q,r)$ because $q$ preserves the set of active blocks,
$\phi_i(ql)=q\phi_i(l)$, and the maps $u_p$ are left $R$-linear.
They also commute with the zero scalar. Thus they define a
$T$-linear map $s:M\to T^N$ with support $\sigma'$. The composite
$fs$ has identity support and acts on every label as $v_pu_p=1$.
Hence $fs=1_M$, proving finite projectivity.

Finally the labels are invariants. A $T$-isomorphism induces a
$Q$-isomorphism of supports and an $R_0$-isomorphism after extension
of scalars along $T\to R_0$. Extension of scalars preserves the
split free presentation and sends its idempotent to $E_0$, so this
last module is exactly $N$. The unique birth quotients in the
$R_0$ classification recover its labels on $D$. Conversely a
support isomorphism and the indicated label isomorphisms give a
fiberwise direct-sum isomorphism commuting with transitions and
scalar action. This proves the stated isomorphism classification.
\end{proof}

For $L=Q$, the Booleanization is $\B$. Taking its single label to be
$R$ recovers the regular free semimodule $T$; taking the label to be
zero gives the projective support semimodule $Q$. If $R=0$, then
$T\cong Q$ and every label is zero, so the theorem reduces exactly
to the original support classification. For $Q=\B$ it is the
classification over $R_0$. No classification of projective
$Q$-semimodules for an arbitrary $Q$ is asserted by these examples.
